% Generated draft: requires mathematical and source review.
\documentclass[11pt]{article}
\usepackage[margin=25mm]{geometry}
\usepackage{fontspec}
\setmainfont{Latin Modern Roman}
\usepackage{amsmath,amssymb,mathtools,unicode-math}
\setmathfont{Latin Modern Math}
\usepackage{enumitem,needspace,xurl,hyperref}
\setlength{\parindent}{0pt}
\setlength{\parskip}{5pt}
\setlength{\emergencystretch}{3em}
\newcommand{\N}{\mathbb{N}}
\newcommand{\Z}{\mathbb{Z}}
\newcommand{\Q}{\mathbb{Q}}
\newcommand{\R}{\mathbb{R}}
\newcommand{\C}{\mathbb{C}}
\newcommand{\F}{\mathbb{F}}
\newcommand{\A}{\mathbb{A}}
\newcommand{\PP}{\mathbb{P}}
\newcommand{\E}{\mathbb{E}}
\newcommand{\eps}{\varepsilon}
\newcommand{\dd}{\,\mathrm{d}}
\begin{document}
\section*{UnsolvedMath ID: 1376}
\textbf{Source problem code:} catalogue entry
\subsection{Definitions and mathematical statement}
Let $(M,g)$ be a smooth Riemannian manifold. For a two-dimensional subspace $\sigma\subset T_pM$, write $K_g(\sigma)$ for its sectional curvature. Strictly positive sectional curvature means $K_g(\sigma)>0$ for every such plane at every point; strictly negative sectional curvature means $K_g(\sigma)<0$ for every such plane. Ricci and scalar curvature impose different conditions. The supplied phrase ``curvature'' does not explicitly select one of these notions, although sectional curvature is the conventional interpretation of many Hopf curvature questions.

When the rational homology groups of $M$ are finite-dimensional and vanish above a finite degree, define its ordinary Euler characteristic by
\[
 \chi(M)=\sum_{i\geq0}(-1)^i b_i(M),\qquad
 b_i(M)=\dim_{\Q}H_i(M;\Q).
\]
A closed manifold means a compact manifold without boundary; this is a familiar setting in which the displayed Euler characteristic is defined. The supplied record does not state closedness, finite homological type, boundary conditions, or whether ordinary rather than compactly supported Euler characteristic is meant.

\textbf{Source research question.} Determine relationships between curvature and Euler characteristic for higher-dimensional Riemannian manifolds. The supplied background gives two purported versions:
\begin{itemize}
\item whether positive curvature on an even-dimensional manifold forces $\chi(M)>0$;
\item whether negative curvature forces $\chi(M)=0$.
\end{itemize}
These are separate imported questions. They are not combined into one theorem, and no unmentioned compactness or curvature hypothesis is silently appended. If positive curvature means positive sectional curvature and $M$ is taken to be closed of dimension $2m$, the first would ask whether
\[
 [K_g(\sigma)>0\text{ for every }\sigma]\ \Longrightarrow\ \chi(M)>0.
\]
This conditional display identifies choices needed for a precise formulation; the source does not explicitly supply them. The negative-curvature/zero-Euler-characteristic wording needs particular verification: it is not interchangeable with an even-dimensional sign assertion involving $(-1)^m\chi(M)$, and no parity or additional hypotheses are attached to it in the record. This draft does not repair that wording by substituting a different Hopf conjecture.

The background motivates these questions through higher-dimensional analogues of Gauss--Bonnet and says special cases have seen progress. The literature triage dated 2026-08-17 emphasizes that ``Hopf Conjectures'' conflates distinct assertions with different statuses. It reports special cases of the Hopf sign conjecture and counterexamples to stronger variants without naming the variants or results, and states that the brief record does not specify a single theorem. Its \emph{partially solved} classification and the registry's \emph{open} label are supplied descriptions of the broad collection, not a verified status for either incomplete formulation above.
\subsection{Short English statement}
The record asks broadly how curvature controls Euler characteristic and mentions distinct positive- and negative-curvature questions. It omits the curvature notion, compactness and boundary assumptions, and exact intended Hopf statement. Its negative-curvature assertion of zero Euler characteristic must be checked rather than silently replaced by a sign conjecture.
\subsection{Sources}
\begin{itemize}
\item UnsolvedMath local registry, data/unsolvedmath/problems.json, permanent ID 1376, source code catalogue entry, ``Hopf Conjectures,'' including the separate background formulations and literature triage dated 2026-08-17.
\item Supplied reference, ``Hopf conjecture overview,'' \url{https://en.wikipedia.org/wiki/Hopf_conjecture}. The review describes this as distinguishing separate curvature conjectures. No particular theorem, primary paper, or authoritative citation supporting the negative-curvature/zero assertion is supplied; the linked overview was not independently consulted here.
\item Ulam AI and the UnsolvedMath Contributors, \emph{UnsolvedMath: A Curated Collection of Open Mathematics Problems}. Dataset curation and original metadata are supplied under CC BY 4.0; referenced source material retains its own terms.
\end{itemize}
\end{document}
